\section*{Chapter \ref{ch:b2:reproductive-hormones} --- Hormonal Control of Reproduction}

\begin{solution}{exo:b2:reproductive-hormones:1}
Hypothalamus (GnRH, pulses) $\to$ anterior pituitary (LH, FSH) $\to$
\omterm{def:b2:mammal-reproduction:testis}{testis}: LH on \omterm{def:b2:mammal-reproduction:testis}{Leydig cells} (testosterone), FSH on \omterm{def:b2:mammal-reproduction:testis}{Sertoli cells}
(\omterm{def:b2:mammal-reproduction:testis}{spermatogenesis}, inhibin). Feedback: testosterone inhibits GnRH and
LH; inhibin inhibits FSH.
\end{solution}

\begin{solution}{exo:b2:reproductive-hormones:2}
Ovary: \omterm{def:b2:reproductive-hormones:cycle}{follicular phase}, days 1--14, oestradiol from the growing
\omterm{def:b2:mammal-reproduction:ovary}{follicle}; \omterm{def:b2:mammal-reproduction:ovary}{ovulation}, day 14, \omterm{def:b2:reproductive-hormones:cycle}{LH surge}; \omterm{def:b2:reproductive-hormones:cycle}{luteal phase}, days 15--28,
progesterone from the \omterm{def:b2:mammal-reproduction:ovary}{corpus luteum}. Uterus: menstruation, days 1--5
(steroid withdrawal); proliferative phase, days 6--14 (oestradiol);
secretory phase, days 15--28 (progesterone).
\end{solution}

\begin{solution}{exo:b2:reproductive-hormones:3}
The daily oestrogen and progestin hold LH and FSH low by negative
feedback, so no \omterm{def:b2:mammal-reproduction:ovary}{follicle} is recruited, no oestradiol rise occurs and
the surge never fires; the progestin also thickens the mucus. In the
pill-free week the steroids fall and the \omterm{def:b2:reproductive-hormones:cycle}{endometrium}, built under
them, is shed --- a withdrawal bleed, not a menstruation after an
\omterm{def:b2:mammal-reproduction:ovary}{ovulation}.
\end{solution}

\begin{solution}{exo:b2:reproductive-hormones:4}
LH and FSH rise several-fold (no testosterone or inhibin feedback);
muscle, libido, the accessory glands regress. Testosterone restores
the characters and brings LH back down, but not fertility: the
testes are gone, and even in a man with testes exogenous
testosterone suppresses LH and so the intratesticular testosterone
that \omterm{def:b2:mammal-reproduction:testis}{spermatogenesis} needs.
\end{solution}

\begin{solution}{exo:b2:reproductive-hormones:5}
$24/1.5 = 16$ pulses a day; $24/4 = 6$. When the \omterm{def:b2:mammal-reproduction:ovary}{corpus luteum} dies
the pulses are still slow, favouring FSH: FSH rises first --- the
right hormone, since it is FSH that recruits the next cohort of
\omterm{def:b2:mammal-reproduction:ovary}{follicles}.
\end{solution}

\begin{solution}{exo:b2:reproductive-hormones:6}
\qty{200}{pg/mL} $= \qty{200}{ng/L}$: $200\times 10^{-9}/272 =
\qty{7.4e-10}{mol/L} = \qty{740}{pmol/L}$. \qty{15}{ng/mL} $=
\qty{15}{\micro g/L}$: $15\times 10^{-6}/314 = \qty{48}{nmol/L}$. In
one microlitre: $7.4\times 10^{-16}$ mol $= 4.4\times 10^{8}$
molecules.
\end{solution}

\begin{solution}{exo:b2:reproductive-hormones:7}
\omterm{def:b2:mammal-reproduction:ovary}{Ovulation} on day $35 - 14 = 21$; fertile from day 18 to day 22. With
irregular cycles the \omterm{def:b2:reproductive-hormones:cycle}{follicular phase}, and so the \omterm{def:b2:mammal-reproduction:ovary}{ovulation} day,
cannot be predicted from past cycles, and the window is missed or
misplaced.
\end{solution}

\begin{solution}{exo:b2:reproductive-hormones:8}
From day 21 to 26 is 2.5 doublings: $2^{2.5} = 5.7$ IU/L --- enough.
Implanting on day 25 gives $1.4$ IU/L on day 26; the \omterm{def:b2:mammal-reproduction:ovary}{corpus luteum},
already regressing, is not rescued, progesterone falls and the
embryo is lost with the period. Late \omterm{def:b2:mammal-reproduction:implantation}{implantation} is a common cause
of early loss.
\end{solution}

\begin{solution}{exo:b2:reproductive-hormones:9}
The pituitary responds to the pattern of GnRH, not its amount:
hourly pulses sustain LH and FSH, the same daily dose infused
continuously abolishes them. It ruled out a simple dose--response.
A month of long-acting agonist: after a few days' flare, LH falls to
near zero (\omterm{prop:b2:reproductive-hormones:pulses}{receptor desensitisation}), oestradiol to postmenopausal
levels, the \omterm{def:b2:reproductive-hormones:cycle}{endometrium} atrophies and no cycle occurs --- a
reversible medical menopause.
\end{solution}

\begin{solution}{exo:b2:reproductive-hormones:10}
No hormone: $R^* = 100/0.5 = 200$. Continuous: $100/6.5 = 15$. A
two-minute pulse removes $6\times 200\times(2/60) = 40$ receptors
(\qty{20}{\%}); in the following hour the deficit shrinks by a
factor $e^{-0.5}$, i.e.\ \qty{39}{\%} of it is recovered. The pool
settles where the loss per pulse equals the recovery per hour:
about 150 receptors before each pulse, three quarters of the
maximum --- ten times the continuous level.
\end{solution}

\begin{solution}{exo:b2:reproductive-hormones:11}
Menopause: no \omterm{def:b2:mammal-reproduction:ovary}{follicles}, so no oestradiol or inhibin, so no
feedback: LH and FSH run high. Prolactinoma: prolactin slows the
GnRH pulse generator, so LH is low, the \omterm{def:b2:mammal-reproduction:ovary}{follicles} are not driven,
oestradiol is low and there is no cycle. Granulosa-cell tumour:
constant high oestradiol suppresses LH by negative feedback (and
without a maturing \omterm{def:b2:mammal-reproduction:ovary}{follicle} there is nothing to ovulate), so the
cycle stops with high oestradiol --- and the \omterm{def:b2:reproductive-hormones:cycle}{endometrium} overgrows.
\end{solution}

\begin{solution}{exo:b2:reproductive-hormones:12}
Low or brief oestradiol reduces GnRH and LH; high, sustained
oestradiol (above \qty{200}{pg/mL} for more than 36 hours) changes
the response of the kisspeptin--GnRH system and of the pituitary,
which has meanwhile stocked LH, so that the same hormone now
provokes a burst. \omterm{def:b2:mammal-reproduction:ovary}{Ovulation} must be all-or-none and timed: the egg
must be released once, on one day, when a \omterm{def:b2:mammal-reproduction:ovary}{follicle} is ripe. A dial
(LH proportional to oestradiol) would give a gradual, partial
luteinisation, \omterm{def:b2:mammal-reproduction:ovary}{follicles} ovulating half-ripe or several at once,
and no defined fertile day.
\end{solution}

\begin{solution}{pb:b2:reproductive-hormones:1}
\textbf{1.} Follicular, days 1--14: oestradiol rising, progesterone
low. Luteal, days 15--28: progesterone high.
\textbf{2.} The surge peaked on day 14 (started on day 13);
\omterm{def:b2:mammal-reproduction:ovary}{ovulation} about 36 hours after the onset, on day 15.
\textbf{3.} Days 15 to 28: 14 days, the lifespan of the \omterm{def:b2:mammal-reproduction:ovary}{corpus
luteum}.
\textbf{4.} Progesterone raises the thermoregulatory set point by
\qty{0.3}{\celsius} to \qty{0.4}{\celsius}; the rise appears a day or
two after \omterm{def:b2:mammal-reproduction:ovary}{ovulation}, so it confirms that \omterm{def:b2:mammal-reproduction:ovary}{ovulation} has happened but
cannot announce it.
\textbf{5.} Days 11 to 14 (210, 260, 310, 220): three to four days,
more than the 36 hours the switch requires.
\textbf{6.} Days 12 to 16.
\textbf{7.} $310\times 10^{-9}/272 = \qty{1.14e-9}{mol/L} =
\qty{1140}{pmol/L}$; $16\times 10^{-6}/314 = \qty{51}{nmol/L}$.
\textbf{8.} $65/8 = 8$. With a one-hour half-life, LH halves every
hour once secretion stops, so a day later it is back near baseline.
\textbf{9.} \omterm{def:b2:mammal-reproduction:ovary}{Ovulation} around day 22, cycle of about 36 days.
\textbf{10.} \omterm{def:b2:mammal-reproduction:implantation}{Implantation} on day 22; by day 27 (five days, 2.5
doublings) about \qty{5.7}{IU/L}.
\textbf{11.} Progesterone falling to \qty{2}{ng/mL} on day 27 means
the \omterm{def:b2:mammal-reproduction:ovary}{corpus luteum} is regressing on schedule: nothing rescued it.
\textbf{12.} Progesterone around \qty{20}{ng/mL} and rising,
oestradiol around \qty{200}{pg/mL}.
\textbf{13.} At the end of the cycle oestradiol, progesterone and
inhibin fall together and FSH escapes its feedback; as the cohort
grows, oestradiol and inhibin rise and FSH falls, and only the
\omterm{def:b2:mammal-reproduction:ovary}{follicle} with the most FSH receptors (the most granulosa cells)
keeps growing on the falling FSH --- the rest die.
\textbf{14.} $R^* = 200$ without hormone; $100/6.5 = 15$ under
continuous hormone.
\textbf{15.} $6\times 200\times 2/60 = 40$ receptors, \qty{20}{\%};
time constant $1/d = \qty{2}{h}$, so \qty{39}{\%} of the deficit is
recovered in an hour; the pool settles around 150 receptors, three
quarters of the maximum.
\textbf{16.} Pulses keep the pituitary at about 150 receptors and it
answers each pulse with LH; a continuous infusion drives the pool to
15 and the same cells fall silent, whatever the dose; pulses restore
the pool and the response.
\textbf{17.} Day one: LH and testosterone rise (the agonist
stimulates receptors that are still there --- the flare). One month:
receptors desensitised, LH near zero, testosterone at castrate
levels, which is the therapeutic aim.
\textbf{18.} FSH rises (no inhibin), LH stays normal; more \omterm{def:b2:mammal-reproduction:ovary}{follicles}
survive selection each cycle, with multiple \omterm{def:b2:mammal-reproduction:ovary}{ovulations} and
dizygotic twins.
\textbf{19.} In a natural cycle the fall of FSH kills all but the
dominant follicle; a constant supply of FSH keeps the whole cohort
growing to maturity.
\textbf{20.} Ten \omterm{def:b2:mammal-reproduction:ovary}{follicles} make far more than \qty{200}{pg/mL} of
oestradiol early, which would trigger a premature \omterm{def:b2:reproductive-hormones:cycle}{LH surge} and
\omterm{def:b2:mammal-reproduction:ovary}{ovulation} before collection; the antagonist blocks GnRH and no
surge can occur.
\textbf{21.} In the table the surge peaked on day 14 and \omterm{def:b2:mammal-reproduction:ovary}{ovulation}
followed about 36 hours after its onset; hCG substitutes for the
surge, and collection at 35 hours catches the oocytes mature but
still in the \omterm{def:b2:mammal-reproduction:ovary}{follicle}.
\textbf{22.} $10\times 0.7 = 7$ embryos, $7\times 0.4 = 2.8$
\omterm{def:b2:mammal-reproduction:implantation}{blastocysts}; $1 - 0.65^{2} = 0.58$.
\textbf{23.} Constant oestrogen and progestin keep FSH and LH low: no
\omterm{def:b2:mammal-reproduction:ovary}{follicle} grows, oestradiol never rises to the threshold, no surge,
no \omterm{def:b2:mammal-reproduction:ovary}{corpus luteum}, no progesterone rise.
\textbf{24.} Testosterone must be given for the secondary characters
and libido; but it does not restore the high intratesticular
concentration that \omterm{def:b2:mammal-reproduction:testis}{Leydig cells} provided, so \omterm{def:b2:mammal-reproduction:testis}{spermatogenesis} stays
suppressed --- which is the point.
\textbf{25.} \omterm{def:b2:mammal-reproduction:ovary}{Ovulation} on day 15; \omterm{def:b2:reproductive-hormones:cycle}{luteal phase} 14 days; fertile
window days 12--16; surge triggered by oestradiol above
\qty{200}{pg/mL} for more than 36 hours, met on days 11--14.
\end{solution}
